% Part: computability
% Chapter: recursive-functions
% Section: normal-form

\documentclass[../../../include/open-logic-section]{subfiles}

\begin{document}

\olfileid{cmp}{rec}{nft}
\olsection{范式定理}

\begin{thm}[Kleene 范式定理]
\ollabel{thm:kleene-nf}
存在原始递归关系 $T(e, x, s)$ 和原始递归函数 $U(s)$，满足以下性质：若 $f$ 为任意部分递归函数，则存在某个~$e$，使得
\[
f(x) \simeq U(\umin{s}{T(e, x, s)})
\]
对所有 $x$ 成立。
\end{thm}

\begin{explain}
范式定理的证明颇为复杂，但其基本思想很简单。
每个部分递归函数都有一个\emph{指标}~$e$，
直观上，它是编码该函数“程序”或定义的数。
若 $f(x) \fdefined$，
则可系统地记录该计算过程并用某个数~$s$ 编码；
而且，仅利用 $x$ 和定义~$e$ 就能原始递归地验证“$s$ 编码了函数~$f$ 在输入~$x$ 上的计算”这一事实。
因此，关系~$T$——即“指标为~$e$ 的函数在输入~$x$ 上有计算，且 $s$ 编码了该计算”——是原始递归的。
给定完整的计算记录~$s$，$s$ 的“最终结果”即为~$f(x)$ 的值，
该值也能从~$s$ 原始递归地提取出来。

范式定理表明，定义任何部分递归函数只需一次无界搜索。
基本上，我们可以搜索所有的数，直到找到一个数，
它编码了指标为~$e$ 的函数在输入~$x$ 上的计算。
我们可以用数~$e$ 作为部分递归函数的“名字”，
并将定理中等式所定义的函数~$f$ 记作 $\cfind{e}$。
注意，任何部分递归函数都可以有多个指标——事实上，每个部分递归函数都有无穷多个指标。
\end{explain}
\end{document}